\documentclass[10pt,a4paper]{amsart}
\usepackage[margin=19mm]{geometry}
\usepackage{amsmath,amssymb,mathtools}
\usepackage[scr=rsfs,cal=euler]{mathalfa}
\usepackage{xfrac}
\usepackage[unicode,hidelinks,bookmarks=false]{hyperref}
\newcommand{\ie}{i.e.\ }
\newcommand{\cf}{cf.\ }
\newcommand{\resp}{respectively\ }
\newcommand{\Subsection}{\subsection}
\providecommand{\nobreakdash}{\nobreak\hbox{-}}
\setlength{\emergencystretch}{2em}
\multlinegap0pt
\usepackage{bm}
\usepackage{bbm}
\usepackage{dsfont}
\usepackage{xspace}
\usepackage{cancel}
\usepackage{mathtools}
\usepackage{amsthm}
\makeatletter
\@ifundefined{thm}{\newtheorem{thm}{Thm}}{}
\@ifundefined{definition}{\newtheorem{definition}[thm]{Definition}}{}
\@ifundefined{lemma}{\newtheorem{lemma}[thm]{Lemma}}{}
\@ifundefined{prop}{\newtheorem{prop}[thm]{Proposition}}{}
\makeatother
\newcommand{\R}{\mathbb R}
\renewcommand{\Re}{\mathop{\rm Re}}
\renewcommand{\bar}{\overline}
\renewcommand{\mod}[1]{{\ifmmode\text{\rm\ (mod~$#1$)}\else\discretionary{}{}{\hbox{\!\!}}\rm(mod~$#1$)\fi}}
\title[Joint distributions of error terms for primes in arithmetic ]{Joint distributions of error terms for primes in arithmetic progressions modulo 11}
\author{K\"ubra Benl\.i, Greg Martin, Paul P\'eringuey}
\date{Source: arXiv:2510.05427v1 (CC BY 4.0)}
\begin{document}
\maketitle

\noindent\emph{Source and licence.} Joint distributions of error terms for primes in arithmetic progressions modulo 11. Version arXiv:2510.05427v1 (CC BY 4.0), distributed under the Creative Commons Attribution 4.0 International licence, as recorded in the arXiv licence metadata for this exact version and shown on the abstract page https://arxiv.org/abs/2510.05427v1. Full rights notice: licenses/arxiv-2510.05427-CC-BY-4.0.txt.\par\smallskip

\noindent\textit{Editorial scope.} Counted unit: the Prop whose source label is \texttt{bound E2} in the article (source0.tex lines 840--842, source label \texttt{bound E2}), delivered together with the article's own complete original proof of it (source0.tex lines 844--876, one \texttt{proof} environment of 33 lines). No mathematics is written, repaired or re-derived: statement and proof are the article's own text line for line. Required context retained uncounted: proposition bound 2 (lines 537--561); labeling characters (lines 705--707); dchi and echi lemma (lines 815--832). Every definition, display and label of those lines is retained unchanged. Omitted from this package and not redistributed: 1--536, 562--704, 708--814, 833--839, 843--843, 877--1017. Nothing inside a retained excerpt is omitted or altered: every retained body is delivered exactly as the article's lines are written, so no omission receipt of any kind accompanies this package. Citations: the retained excerpts carry no citation command at all, so no bibliography entry of the article is transcribed and no reference is redistributed. Reference adaptation: none was needed and none was made; every reference inside the delivered unit resolves inside it, so no unresolved reference or citation is produced. Printed slips kept literal: no printed word, symbol, hypothesis or number of the counted statement or of its proof was altered. Layout and class: the article's own class is \texttt{article} with packages including fullpage, babel, amsthm, amsfonts, amssymb, amsmath, mathrsfs, textcomp, amscd, amsbsy, hyperref, enumitem, eurosym, bbm; the wrapper is the lane's pinned \texttt{amsart} editorial wrapper at 10pt on A4, which restates the theorem-like environments the retained text needs and adds the article's own macro definitions that the retained text needs (\texttt{\textbackslash R}, \texttt{\textbackslash Re}, \texttt{\textbackslash bar}, \texttt{\textbackslash mod}), with extra wrapper packages (bm, bbm, dsfont, xspace, cancel, mathtools). The delivered numbering is the unit's own. Assets: no retained span contains a figure environment, a graphics-inclusion command, a PDF-page inclusion, a table or a TikZ picture; no image file is copied into this package, the package declares no asset dependency and no image file is redistributed with it. Licence: version arXiv:2510.05427v1 of this article is distributed under the Creative Commons Attribution 4.0 International licence, as recorded in the arXiv licence metadata for this exact version and shown on the abstract page https://arxiv.org/abs/2510.05427v1. A full rights notice is delivered at licenses/arxiv-2510.05427-CC-BY-4.0.txt. Characters: every retained excerpt is pure ASCII, so the delivered engine, XeLaTeX with its default Latin Modern text font, reports no missing character.
\begin{prop}\label{proposition bound 2}
Fix a modulus $q\ge3$. For each nontrivial character $\chi\mod q$, fix positive constants $d_+(\chi),d_-(\chi),d(\chi)$ and $e_+(\chi),e_-(\chi),e(\chi)$ such that
\begin{align} \label{F bound polynomial}
|F(x,\chi)|\le \min\{ 1, d_+(\chi)|x|^{-e_+(\chi)}, d_-(\chi)|x|^{-e_-(\chi)}, d(\chi)|x|^{-e(\chi)} \}
\end{align}
for all $x\in\R$.
Choose real numbers $b>1$, $c\in[0,1]$, $c_+\in[0,1]$ and $c_-\in[-1,0]$, and set
\[
\begin{alignedat}{2}
d_{c_+} &= \prod_{\substack{\chi \mod{q}\\ \Re(\chi(a))\ge c_+}} d_+(\chi) &\quad\text{and}\quad& e_{c_+} = \sum_{\substack{\chi \mod{q}\\ \Re(\chi(a))\ge c_+}}e_+(\chi) \\
d_{c_-} &= \prod_{\substack{\chi \mod{q}\\ \Re(\chi(a))\le c_-}} d_-(\chi) &\quad\text{and}\quad& e_{c_-} = \sum_{\substack{\chi \mod{q}\\ \Re(\chi(a))\le c_-}} e_-(\chi) \\
d_c &= \prod_{\substack{\chi \mod{q}\\ |\Re(\chi(a))|\le c}} d(\chi) &\quad\text{and}\quad& e_c = \sum_{\substack{\chi \mod{q}\\ |\Re(\chi(a))|\le c}} e(\chi).
\end{alignedat}
\]
If $e_c>1$, then for any real numbers $\varepsilon>0$ and $C\ge1$,
\begin{equation} \label{triple threat}
B_2(a,\varepsilon,C)\le B_2^+(a,\varepsilon,C)+B_2^-(a,\varepsilon,C)+\widetilde{B}_2(a,\varepsilon,C),
\end{equation}
where
\begin{align*}
B_2^+(a,\varepsilon,C) &= \frac{bd_{c_+}}{4e_{c_+}}\left(1-\frac1{b}\right)\left(\varepsilon\left(1 +\frac{c_+}{b}\right)\right)^{-e_{c_+}}\left(\left\lfloor\frac{C}{2}\right\rfloor-1\right)^{-e_{c_+}} \\
B_2^-(a,\varepsilon,C) &= \frac{bd_{c_-}}{4e_{c_-}}\left(1-\frac1{b}\right)\left(\varepsilon\left(1 -\frac{c_-}{b}\right)\right)^{-e_{c_-}}\left(\left\lfloor\frac{C}{2}\right\rfloor-1\right)^{-e_{c_-}} \\
\widetilde{B}_2(a,\varepsilon,C) &= \frac{d_c}{e_c-1}\left(\frac C{2b}+1 \right) \left(\varepsilon\left(1-\frac{c}{b}\right)\right)^{-e_c}\left(\left\lfloor\frac{C}{2}\right\rfloor-1\right)^{-e_c}.
\end{align*}
\end{prop}

\begin{definition} \label{labeling characters}
For $0\le j\le 9$ we write $\chi_j$ for the Dirichlet character$\mod{11}$ such that $\chi_j(2) = (e^{2\pi i/10})^j$ (so that $\chi_0$ is the principal character as usual). Note that in this labeling, $\bar\chi_j = \chi_{10-j}$ for $1\le j\le 9$, and in particular $\chi_5$ is the quadratic character (the Legendre symbol) modulo~$11$. This labeling is the same as the labeling in the SageMath library.
\end{definition}

\begin{lemma} \label{dchi and echi lemma}
Assume GRH. Let $\chi_1,\dots,\chi_9$ be the nonprincipal characters modulo $q=11$, labeled as in Definition~\ref{labeling characters}. For each pair $(d(\chi),e(\chi))$ given in the table below, we have $|F(x,\chi)| \le \min\{1,d(\chi)|x|^{-e(\chi)} \}$ for all $x\in\R$.
\begin{center}$
\begin{array}{|c|r|r|r|r|}
\hline
\chi & \substack{\textstyle d(\chi)\text{ when} \\ \textstyle e(\chi)=5} & \substack{\textstyle d(\chi)\text{ when} \\ \textstyle e(\chi)=8.5} & \substack{\textstyle d(\chi)\text{ when} \\ \textstyle e(\chi)=9.5} & \substack{\textstyle d(\chi)\text{ when} \\ \textstyle e(\chi)=10} \\ \hline
\chi_1 & 820 & 1{,}855{,}630 & 21{,}021{,}079 & 73{,}516{,}699 \\ \hline
\chi_2 & 1{,}189 & 2{,}875{,}162 & 32{,}845{,}058 & 115{,}861{,}968 \\ \hline
\chi_3 & 1{,}195 & 2{,}916{,}371 & 33{,}450{,}847 & 116{,}963{,}421 \\ \hline
\chi_4 & 1{,}203 & 2{,}950{,}376 & 34{,}133{,}529 & 119{,}474{,}311 \\ \hline
\chi_5 & 678 & 1{,}549{,}125 & 17{,}785{,}195 & 61{,}586{,}977 \\ \hline
\chi_6 & 770 & 1{,}800{,}290 & 20{,}607{,}026 & 71{,}566{,}480 \\ \hline
\chi_7 & 355 & 742{,}204 & 8{,}398{,}441 & 28{,}913{,}287 \\ \hline
\chi_8 & 880 & 2{,}026{,}172 & 23{,}375{,}053 & 81{,}872{,}319 \\ \hline
\chi_9 & 715 & 1{,}590{,}237 & 18{,}149{,}703 & 63{,}453{,}141 \\ \hline
\end{array}
$\end{center}
\end{lemma}

\begin{prop} \label{bound E2}
Assume GRH. For $q=11$ and any $2\le a\le 10$, we have $|E_2(0.2,100)| \le 2.44\times 10^{-11}$.
\end{prop}

\begin{proof}
In Proposition~\ref{proposition bound 2},
we take $b=4$ and $c=1$ and $e(\chi)=5$.
For each $2\le a\le10$,
we carefully manually choose values for $c_+$, $c_-$, $e_+(\chi)$, and $e_-(\chi)$. The values for $d_+(\chi),d_-(\chi),d(\chi)$ are then determined from $e_+(\chi),e_-(\chi),e(\chi)$
by Lemma~\ref{dchi and echi lemma}. Our choices, and the corresponding conclusions from Proposition~\ref{proposition bound 2}, are recorded in the following table.
\begin{center}$
\begin{array}{|c | c c c c | c |} 
\hline
a & c_+ & c_- & e_+(\chi) & e_-(\chi) & B_2(a,0.2,100)\le \\
\hline\hline
2 & 0.309 & -0.309 & 8.5 & 8.5 & 6.07\times 10^{-12} \\ 
\hline
3 & 0.309 & -0.809 & 8.5 & 9.5 & 1.50\times 10^{-13} \\ 
\hline
4 & 0.309 & -0.809 & 8.5 & 9.5 & 1.50\times 10^{-13} \\ 
\hline
5 & 0.309 & -0.809 & 8.5 & 9.5 & 1.67\times 10^{-13} \\ 
\hline
6 & 0.309 & -0.309 & 8.5 & 8.5 & 6.07\times 10^{-12} \\ 
\hline
7 & 0.309 & -0.309 & 8.5 & 8.5 & 4.09\times 10^{-12} \\ 
\hline
8 & 0.309 & -0.309 & 8.5 & 8.5 & 4.09\times 10^{-12} \\ 
\hline
9 & 0.309 & -0.809 & 8.5 & 9.5 & 1.67\times 10^{-13} \\ 
\hline
10 & 1 & -1 & 10 & 10 & 9.72\times 10^{-14} \\
\hline
\end{array}
$\end{center}
In all cases, $|E_2(0.2,100)|\le 4B_2(a,0.2,100) \le 2.44\times 10^{-11}$ as claimed.
\end{proof}

\end{document}
